\chapter{Explanatory Depth}\label{ch:depth}

A second natural dimension of variation is how much a film explains. A documentary about a proof can pause to define its terms or assume them. An instructional film can show every step of a repair or compress the familiar ones. A historical drama can supply background or trust the audience to have it. Realizations that differ in this way are usually called beginner and expert versions, and this chapter uses those names while arguing that both are misleading.

\section{What differs}

Beginner and expert realizations do not differ in intelligence, and neither is a simplified form of the other. They differ in what they presume. The expert realization presumes that its viewers bring certain background knowledge; the beginner realization supplies that background on screen. Both can be equally demanding. The beginner realization asks its viewers to learn something during the film; the expert realization asks them to build on something learned before it.

In the terms of \Cref{ch:switch-admissibility}, the difference is in how each realization comes to hold its background. Call the background a realization needs $B$. The expert realization \emph{declares} $B$: it treats it as established from the first frame. The beginner realization \emph{establishes} $B$ gradually, explanation by explanation, so that at time $t$ it has established some part $B(t)\subseteq B$, growing until, at some time $t^{*}$, it has established all of it.

\section{The one-way window}

This structure produces the asymmetry noted in \Cref{ch:divergence}, and makes it exact.

Before $t^{*}$, a viewer of the expert realization holds all of $B$, by declaration, and the beginner continuation presupposes only $B(t)$. The switch from expert to beginner satisfies coverage. A viewer of the beginner realization holds only $B(t)$, and the expert continuation may presuppose any part of $B$. The switch from beginner to expert fails coverage whenever the expert continuation relies on background not yet explained. After $t^{*}$, the beginner viewer holds all of $B$, and the switch becomes possible in both directions, other things being equal.

The interval before $t^{*}$ is therefore a \emph{one-way window}. Transitions in the switch system run from the realization that assumes more toward the one that explains more, and $t^{*}$, the moment the explanation is complete, is where they become two-way. Writers of a depth family can place $t^{*}$ deliberately. A family that front-loads explanation reaches it early, and its realizations become mutually accessible for most of the film. A family that explains as it goes reaches it late, or never, and keeps the window open throughout.

The asymmetry has a practical meaning. A viewer who chose the expert realization and finds a passage too hard can move into the beginner realization during the window and follow the explanation there. A viewer who chose the beginner realization cannot jump ahead into the expert one until the explanation is done. Moving toward more help is always possible; moving toward less help must wait for the help to be complete. That is, on reflection, how learning works.

\section{A presumption about people}

The formal account has a weak point, and it should be named rather than hidden. Counting the expert realization's declared background as ``established'' is a presumption about its viewers. The film has not shown them $B$. It assumes they already have it. The switch system inherits the assumption: it treats any viewer on the expert realization as holding $B$, whether or not they do.

This matters in two ways. First, it means the switch system cannot protect a viewer who chose the expert realization without the background. That viewer will be lost, and the system will report that their history is coherent. Second, it means that whoever decides which viewers are on which realization is making a judgment about what those viewers know. If viewers choose for themselves, the judgment is theirs, and a wrong choice is a mistake they can correct by moving toward the beginner realization, which the one-way window always permits. If someone else assigns them, the judgment is made about them, and it can be wrong in ways that are not merely inconvenient. A student placed on the beginner realization by a teacher is being told something about themselves. \Cref{ch:instruction} returns to this. The principle this chapter can state is that depth should be a property of the presentation a viewer selects, not a rank assigned to the viewer.

\section{Where the time comes from}

Explanation takes time. A beginner realization that pauses to define a term, show an intermediate step, or supply background runs longer than an expert realization that does not, unless the time is found somewhere. In the vocabulary of \Cref{ch:family}, explanation spends time distance, and the envelopes of \Cref{ch:genre} require it to be repaid before each scene ends. There are four places to find it.

The first is redundancy. Films contain time that carries little information: establishing shots, transitions, lingering reaction shots. A beginner realization can replace some of it with explanation. The cost falls on atmosphere, which the beginner realization may miss.

The second is to use the same minutes differently. An expert realization need not be a beginner realization with the explanations cut out, a shorter film padded to length. It can spend the time the beginner realization uses for explanation on something the expert viewer will value: a further implication, a harder example, a historical aside, a complication the beginner realization cannot afford. The two realizations then have equal length and different content, and each is complete on its own terms. This is the strongest arrangement, because it makes the expert realization a different film rather than a lesser one, and it keeps the two synchronized without either waiting for the other.

The third is sound. Much explanation is verbal, and verbal explanation can in principle be carried by a separate audio channel while the picture stays common. This requires individual sound, and \Cref{ch:composite-sound} shows that individual sound has costs of its own. If the family keeps a shared score, explanation must be carried by the picture: diagrams, captions, demonstrations, visual emphasis. That is a real constraint, and it shapes what a depth family can be. A shared-sound depth family is closer to a well-designed diagram than to a lecture.

The fourth is flexibility at the boundaries. A family can leave some intervals before its rendezvous deliberately loose, with material in the expert realization that can be shortened or lengthened to absorb whatever time the beginner realization needs. This is the least elegant solution and often the most practical.

\section{Presentation is not uptake}

There is a deeper limit, and it concerns what ``established'' means. The model of \Cref{ch:switch-admissibility} counts a proposition as established when the film presents it as settled. But presentation is not uptake. A fact shown in the corner of a frame, or mentioned once in a line of dialogue, is presented to every viewer and taken up by some. Depth variation is often less about what is presented than about whether it is taken up: the beginner realization shows the same thing more slowly, more prominently, more than once.

Mystery films make the point sharply. A beginner realization of a mystery might foreground its clues, holding on them, returning to them, so that a viewer is likely to notice. An expert realization might present the same clues in passing, leaving the viewer to catch them. In the literal model the two realizations establish the same facts, and the realization graph is complete. In practice they differ. A viewer who watched the expert realization throughout and missed a clue is, by the model, fully prepared for its solution, and may be baffled by it. A viewer who switches into the expert realization from the beginner one at the solution arrives better prepared than the viewers who were there all along.

The model can be extended to track salience, weighting what is established by how prominently it was shown, but the extension is uncertain, because uptake varies between people. It is better to state the limit plainly. Switch admissibility guarantees that a viewer was shown what they need. It cannot guarantee that they took it in. For most variation this gap is small. For depth variation it is the whole subject, and the switch system should be read as a floor on coherence, not a promise of comprehension.

\section{Depth as a gift}

The chapter's argument can be summarized by what depth variation offers when it is done well. Every viewer in the room follows the same film, in step, at the level they chose. Those who want more explanation receive it without slowing anyone else, and those who want less receive something in its place rather than a pause. Any viewer who finds their choice too hard can move toward more help at any moment, and anyone who chose more help can move toward less once the help is complete. Nobody is ranked. That is a considerable thing to be able to offer in a single screening, and it is available only because the explanation happens in time a common screen would otherwise have to share.
